% Exercitium2018.tex 
% Exercitium (curso 2018-19)
% José A. Alonso Jiménez
% Sevilla, 2 de febrero de 2019
% ======================================================================

\documentclass[a4paper,12pt,twoside]{book}

%%%%%%%%%%%%%%%%%%%%%%%%%%%%%%%%%%%%%%%%%%%%%%%%%%%%%%%%%%%%%%%%%%%%%%%%
%% § Paquetes adicionales
%%%%%%%%%%%%%%%%%%%%%%%%%%%%%%%%%%%%%%%%%%%%%%%%%%%%%%%%%%%%%%%%%%%%%%%%

% Configuración para XeLaTeX
\usepackage{fontspec}
\usepackage{xltxtra}
\defaultfontfeatures{Ligatures=TeX,Numbers=OldStyle}
\setromanfont{DejaVu Sans}
% \setsansfont{Arial}
\setmonofont{DejaVu Sans Mono}[Scale={0.90}]

% Notas: La lista de fuentes disponibles se obtiene con fc-list

% \usepackage{ucs}
% \usepackage[utf8]{inputenc}        % Acentos de UTF8
\usepackage[spanish]{babel}        % Castellanización.
% \usepackage[T1]{fontenc}           % Codificación T1 con European Computer
%                                    % Modern.  
% \usepackage{graphicx}
% \usepackage{fancyvrb}              % Verbatim extendido
% \usepackage{mathpazo}              % Fuentes semejante a palatino
% \usepackage[scaled=.90]{helvet}
% \usepackage{cmtt}
% \renewcommand{\ttdefault}{cmtt}
\usepackage{a4wide}
\usepackage{minted}
\usepackage{comment}
\usepackage{amssymb, amsmath, amsbsy}

\usepackage{titletoc}
\dottedcontents{chapter}[0em]{}{32em}{1pc}
% \dottedcontents{chapter}[<left>]{<above-code>}{<label width>}{<leader width>}

\linespread{1.05}                  % Distancia entre líneas
\setlength{\parindent}{2em}        % Indentación de comienzo de párrafo
\setlength{\parskip}{1ex}          % Distancia entre párrafos
% \deactivatetilden                  % Elima uso de ~ para la eñe
\raggedbottom                      % No ajusta los espacios verticales.

\usepackage[%
  colorlinks=true,
  urlcolor=blue,
  % pdftex,
  pdfauthor={José A. Alonso <jalonso@us.es>},%
  pdftitle={Exercitium (curso 2018-19)},%
  pdfstartview=FitH,%
  bookmarks=false]{hyperref}      

\setcounter{tocdepth}{1}
\setcounter{secnumdepth}{4}

\usepackage{tocstyle}
\usetocstyle{KOMAlike}

% \usepackage{tocloft}
% \renewcommand\cftpartnumwidth{3cm}

% \usepackage{minitoc}

% \setlength\cftparskip{-2pt}
% \setlength\cftbeforechapskip{0pt}

%%%%%%%%%%%%%%%%%%%%%%%%%%%%%%%%%%%%%%%%%%%%%%%%%%%%%%%%%%%%%%%%%%%%%%%%
%% § epigraph: Para citas al principio del capítulo                   %%
%%%%%%%%%%%%%%%%%%%%%%%%%%%%%%%%%%%%%%%%%%%%%%%%%%%%%%%%%%%%%%%%%%%%%%%%

% \usepackage{epigraph} % Para citas al principio del capítulo
\usepackage{epigraph,varwidth}

\renewcommand{\epigraphsize}{\small}
\setlength{\epigraphwidth}{0.6\textwidth}
\renewcommand{\textflush}{flushleft}
\renewcommand{\sourceflush}{flushleft}
% A useful addition
\newcommand{\epitextfont}{\itshape}
% \newcommand{\episourcefont}{\scshape}
\newcommand{\episourcefont}{\normalfont}

\makeatletter
\newsavebox{\epi@textbox}
\newsavebox{\epi@sourcebox}
\newlength\epi@finalwidth
\renewcommand{\epigraph}[2]{%
  \vspace{\beforeepigraphskip}
  {\epigraphsize\begin{\epigraphflush}
   \epi@finalwidth=\z@
   \sbox\epi@textbox{%
     \varwidth{\epigraphwidth}
     \begin{\textflush}\epitextfont#1\end{\textflush}
     \endvarwidth
   }%
   \epi@finalwidth=\wd\epi@textbox
   \sbox\epi@sourcebox{%
     \varwidth{\epigraphwidth}
     \begin{\sourceflush}\episourcefont#2\end{\sourceflush}%
     \endvarwidth
   }%
   \ifdim\wd\epi@sourcebox>\epi@finalwidth 
     \epi@finalwidth=\wd\epi@sourcebox
   \fi
   \leavevmode\vbox{
     \hb@xt@\epi@finalwidth{\hfil\box\epi@textbox}
     \vskip1.75ex
     \hrule height \epigraphrule
     \vskip.75ex
     \hb@xt@\epi@finalwidth{\hfil\box\epi@sourcebox}
   }%
   \end{\epigraphflush}
   \vspace{\afterepigraphskip}}}
\makeatother

%%%%%%%%%%%%%%%%%%%%%%%%%%%%%%%%%%%%%%%%%%%%%%%%%%%%%%%%%%%%%%%%%%%%%%%%%%%%%%
%% § Cabeceras                                                              %%
%%%%%%%%%%%%%%%%%%%%%%%%%%%%%%%%%%%%%%%%%%%%%%%%%%%%%%%%%%%%%%%%%%%%%%%%%%%%%%

\usepackage{fancyhdr}

\addtolength{\headheight}{\baselineskip}

\pagestyle{fancy}

\cfoot{}

\fancyhead{}
\fancyhead[RE]{\mdseries\sffamily Exercitium (2018--19)}
\fancyhead[LO]{\mdseries\sffamily \nouppercase{\leftmark}}
\fancyhead[LE,RO]{\mdseries\sffamily \thepage}

%%%%%%%%%%%%%%%%%%%%%%%%%%%%%%%%%%%%%%%%%%%%%%%%%%%%%%%%%%%%%%%%%%%%%%%%
%% § Definiciones
%%%%%%%%%%%%%%%%%%%%%%%%%%%%%%%%%%%%%%%%%%%%%%%%%%%%%%%%%%%%%%%%%%%%%%%%

\input definiciones
\def\mtctitle{Contenido}

%%%%%%%%%%%%%%%%%%%%%%%%%%%%%%%%%%%%%%%%%%%%%%%%%%%%%%%%%%%%%%%%%%%%%%%%
%% § Título
%%%%%%%%%%%%%%%%%%%%%%%%%%%%%%%%%%%%%%%%%%%%%%%%%%%%%%%%%%%%%%%%%%%%%%%%

\title{
  {\LARGE Exercitium (curso 2018--19) \\
  {\Large Ejercicios de programación funcional con Haskell \\
  {\normalsize (hasta el 25 de enero de 2019)}}} }  
\author{\href{http://www.cs.us.es/~jalonso}
        {\Large José A. Alonso Jiménez}}
\date{\vfill \hrule \vspace*{2mm}
  \begin{tabular}{l}
      \href{http://www.cs.us.es/glc}
           {Grupo de Lógica Computacional} \\
      \href{http://www.cs.us.es}
           {Dpto. de Ciencias de la Computación e Inteligencia Artificial} \\
      \href{http://www.us.es}
           {Universidad de Sevilla}  \\
      Sevilla, \today 
  \end{tabular}\hfill\mbox{}}

%%%%%%%%%%%%%%%%%%%%%%%%%%%%%%%%%%%%%%%%%%%%%%%%%%%%%%%%%%%%%%%%%%%%%%%%
%% § Documento
%%%%%%%%%%%%%%%%%%%%%%%%%%%%%%%%%%%%%%%%%%%%%%%%%%%%%%%%%%%%%%%%%%%%%%%%

% \includeonly{Números con dígitos 1 y 2}

% \includexmp{licencia}

\begin{document}
% \dominitoc

\maketitle
\newpage

\noindent
Esta obra está bajo una licencia Reconocimiento--NoComercial--CompartirIgual
2.5 Spain de Creative Commons. 

\vspace*{4ex}

\begin{center}
\fbox{
\begin{minipage}{\textwidth}
\vspace*{1ex}
\noindent                     
{\bf Se permite:}
\begin{itemize}                         
\item copiar, distribuir y comunicar públicamente la obra
\item hacer obras derivadas
\end{itemize}
         
\noindent       
{\bf Bajo las condiciones siguientes:}

\begin{minipage}[t]{20mm}
  \vspace*{0mm}\hspace*{1em}
  \includegraphics[scale=0.35]{deed.png}
\end{minipage}%
\begin{minipage}[t]{14cm}
  \vspace*{0mm}
  {\bf Reconocimiento}. Debe reconocer los créditos de la  obra de la manera
  especificada por el autor.  
\end{minipage}
\\[1ex]

\begin{minipage}[t]{20mm}
  \vspace*{0mm}\hspace*{1em}
  \includegraphics[scale=0.35]{deed-eu.png}
\end{minipage}%
\begin{minipage}[t]{14cm}
  \vspace*{2mm}
  {\bf No comercial}. No puede utilizar esta obra para fines comerciales.
\end{minipage}
\\[1ex]

\begin{minipage}[t]{20mm}
  \vspace*{0mm}\hspace*{1em}
  \includegraphics[scale=0.35]{deed_002.png}
\end{minipage}%
\begin{minipage}[t]{14cm}
  \vspace*{-2mm}
  {\bf Compartir bajo la misma licencia}. Si altera o transforma esta obra, o
  genera una obra derivada, sólo puede distribuir la obra generada bajo una
  licencia idéntica a ésta. 
\end{minipage}

\begin{itemize}    
\item Al reutilizar o distribuir la obra, tiene que dejar bien claro los
  términos de la licencia de esta obra. 
\item Alguna de estas condiciones puede no aplicarse si se obtiene el permiso
 del titular de los derechos de autor.
\end{itemize}
\vspace*{1ex}
\end{minipage}}
\end{center}

\vspace*{4ex}

Esto es un resumen del texto legal (la licencia completa). Para ver una copia
de esta licencia, visite 
\href{http://creativecommons.org/licenses/by-nc-sa/2.5/es/}
     {\url{http://creativecommons.org/licenses/by-nc-sa/2.5/es/}}
o envie una carta a Creative Commons, 559 Nathan Abbott Way, Stanford,
California 94305, USA.

%%% Local Variables: 
%%% mode: latex
%%% TeX-master: t
%%% End: 

\newpage

\newpage

\mbox{} \vspace*{2cm}
  \begin{verse}
  ``Sorpresas tiene la vida, \\
  Guiomar, del alma y del cuerpo; \\ 
  que nadie guarde hasta el fin \\
  el nombre que le pusieron; \\
  nadie crea ser quien dicen \\
  que es, ni que pueda serlo.'' \\ \vspace*{2ex}

  De Antonio Machado
  \end{verse}

\begin{flushright} 
\textit{Para Guiomar}
\end{flushright}

\newpage

\tableofcontents
\clearpage

\renewcommand{\chaptername}{Ejercicio}

\chapter*{Introducción}

% \mbox{} \hspace*{1cm} 

\begin{quote}
  ``\textit{The chief goal of my work as an educator and author is to
  help people learn to write beautiful programs.}''

  (Donald Knuth en
  \href{http://www.paulgraham.com/knuth.html}{Computer programming as an art})
\end{quote}

\vspace* {1cm}

Este libro es una recopilación de las soluciones de los ejercicios
propuestos en el blog
\href{https://www.glc.us.es/~jalonso/exercitium}
     {Exercitium}\
     \footnote{\url{https://www.glc.us.es/~jalonso/exercitium}}
durante el curso 2018--19.

El principal objetivo de Exercitium es servir de complemento a la
asignatura de
\href{https://www.cs.us.es/~jalonso/cursos/i1m-18}
     {Informática}\
     \footnote{\url{https://www.cs.us.es/~jalonso/cursos/i1m-18}}
de 1º del Grado en Matemáticas de la Universidad de Sevilla.

Con los problemas de Exercitium, a diferencias de los de las
\href{https://www.cs.us.es/~jalonso/cursos/i1m-18/ejercicios/ejercicios-I1M-2018.pdf}
     {relaciones}\
     \footnote{\url{https://www.cs.us.es/~jalonso/cursos/i1m-18/ejercicios/ejercicios-I1M-2018.pdf}},
se pretende practicar con los conocimientos adquiridos durante todo el
curso, mientras que con las relaciones están orientadas a los nuevos
conocimientos.

Habitualmente de cada ejercicio se muestra distintas soluciones y se
compara sus eficiencias.

La dinámica del blog es la siguiente: cada día, de lunes a viernes, se
propone un ejercicio para que los alumnos escriban distintas soluciones
en los comentarios. Pasado 7 días de la propuesta de cada ejercicio, se
cierra los comentarios y se publica una selección de sus soluciones.

Para conocer la cronología de los temas explicados se puede consultar el
\href{https://www.glc.us.es/~jalonso/vestigium/category/curso/i1m/i1m2018}
     {diario de clase}\
     \footnote{\url{https://www.glc.us.es/~jalonso/vestigium/category/curso/i1m/i1m2018}}.

En el libro se irán añadiendo semanalmente las soluciones de los
ejercicios del curso.

El código del libro se encuentra en
\href{https://github.com/jaalonso/Exercitium2018}
     {GitHub}\
     \footnote{\url{https://github.com/jaalonso/Exercitium2018}}

\section*{Cuaderno de bitácora}

En esta sección se registran los cambios realizados en las sucesivas
versiones del libro.

\subsection*{Versión del 16 de diciembre de 2018}

Se han añadido los ejercicios resueltos de la primera semana de
diciembre:

\begin{itemize}
\item \nameref{031218}
\item \nameref{041218}
\item \nameref{051218}
\item \nameref{061218}
\item \nameref{071218}
\end{itemize}
     
\subsection*{Versión del 22 de diciembre de 2018}

Se han añadido los ejercicios resueltos de la primera semana de
diciembre:

\begin{itemize}
\item \nameref{101218}
\item \nameref{111218}
\item \nameref{121218}
\item \nameref{131218}
\item \nameref{141218}
\end{itemize}
     
\subsection*{Versión del 29 de diciembre de 2018}

Se han añadido los ejercicios resueltos de la primera semana de
diciembre:

\begin{itemize}
\item \nameref{171218}
\item \nameref{181218}
\item \nameref{191218}
\item \nameref{201218}
\item \nameref{211218}
\end{itemize}
     
\subsection*{Versión del 29 de diciembre de 2018}

Se han añadido los ejercicios resueltos del 24 al 28 de
diciembre:

\begin{itemize}
\item \nameref{181224}
\item \nameref{181225}
\item \nameref{181226}
\item \nameref{181227}
\item \nameref{181228}
\end{itemize}

\subsection*{Versión del 12 de enero de 2019}

Se han añadido los ejercicios resueltos del 31 de diciembre al 4 de
enero: 

\begin{itemize}
\item \nameref{181231}
\item \nameref{190101}
\item \nameref{190102}
\item \nameref{190103}
\item \nameref{190104}
\end{itemize}

\subsection*{Versión del 19 de enero de 2019}

Se han añadido los ejercicios resueltos hasta el 11 de enero:

\begin{itemize}
\item \nameref{190107}
\item \nameref{190108}
\item \nameref{190109}
\item \nameref{190110}
\item \nameref{190111}
\end{itemize}

\subsection*{Versión del 26 de enero de 2019}

Se han añadido los ejercicios resueltos hasta del 14 al 18 de enero:

\begin{itemize}
\item \nameref{190114}
\item \nameref{190115}
\item \nameref{190116}
\item \nameref{190117}
\item \nameref{190118}
\end{itemize}

\subsection*{Versión del 2 de febrero de 2019}

Se han añadido los ejercicios resueltos hasta del 21 al 25 de enero:

\begin{itemize}
\item \nameref{190121}
\item \nameref{190122}
\item \nameref{190123}
\item \nameref{190124}
\item \nameref{190125}
\end{itemize}

\chapter{Listas equidigitales}
\entrada{Listas_equidigitales}

\chapter{Distancia de Hamming}
\entrada{Distancia_de_Hamming}

\chapter{Último dígito no nulo del factorial}
\entrada{Ultimo_digito_no_nulo_del_factorial}

\chapter{Diferencia simétrica}
\entrada{Diferencia_simetrica}

\chapter{Números libres de cuadrados}
\entrada{Numeros_libres_de_cuadrados}

\chapter{Capicúas productos de dos números de dos dígitos}
\entrada{Capicuas_productos_de_dos_numeros_de_dos_digitos}

% ---------------------------------------------------------------------
 
\chapter{Números autodescriptivos}
\entrada{Numeros_autodescriptivos}

\chapter{Número de parejas}
\entrada{Numeros_de_parejas}

\chapter{Reconocimiento de particiones}
\entrada{Reconocimiento_de_particiones}

\chapter{Relación definida por una partición}
\entrada{Relacion_definida_por_una_particion}

\chapter{Ceros finales del factorial}
\entrada{Ceros_finales_del_factorial}

% ---------------------------------------------------------------------
 
\chapter{Números primos sumas de dos primos}
\entrada{Numeros_primos_sumas_de_dos_primos}

\chapter{Suma de inversos de potencias de cuatro}
\entrada{Suma_de_inversos_de_potencias_de_cuatro}

\chapter{Elemento solitario}
\entrada{Elemento_solitario}

\chapter{Números colinas}
\entrada{Numeros_colinas}

\chapter{Raíz cúbica entera}
\entrada{Raiz_cubica_entera}

% ---------------------------------------------------------------------

\chapter{Numeración de los árboles binarios completos}
\label{031218}
\entrada{Numeracion_de_arboles_binarios_completos}

\chapter{Posiciones en árboles binarios}
\label{041218}
\entrada{Posiciones_en_arboles_binarios}

\chapter{Posiciones en árboles binarios completos}
\label{051218}
\entrada{Posiciones_en_arboles_binarios_completos}

\chapter{Elemento del árbol binario completo según su
  posición}
\label{061218}
\entrada{Elemento_del_arbol_binario_completo_segun_su_posicion}

\chapter{Aproximación entre pi y e}
\label{071218}
\entrada{Aproximacion_entre_pi_y_e}

% ---------------------------------------------------------------------

\chapter{Menor contenedor de primos}
\label{101218}
\entrada{Menor_contenedor_de_primos}

\chapter{Árbol de computación de Fibonacci}
\label{111218}
\entrada{Arbol_de_computacion_de_Fibonacci}

\chapter{Entre dos conjuntos}
\label{121218}
\entrada{Entre_dos_conjuntos}

\chapter{Expresiones aritméticas generales}
\label{131218}
\entrada{Expresiones_aritméticas_generales}

\chapter{Superación de límites}
\label{141218}
\entrada{Superacion_de_limites}

% ---------------------------------------------------------------------

\chapter{Intercambio de la primera y última columna de
  una matriz}
\label{171218}
\entrada{Intercambio_de_la_primera_y_ultima_columna_de_una_matriz}

\chapter{Números primos de Pierpont}
\label{181218}
\entrada{Numeros_primos_de_Pierpont}

\chapter{Grado exponencial}
\label{191218}
\entrada{Grado_exponencial}

\chapter{Divisores propios maximales}
\label{201218}
\entrada{Divisores_propios_maximales}

\chapter{Árbol de divisores}
\label{211218}
\entrada{Arbol_de_divisores}

% ---------------------------------------------------------------------

\chapter{Divisores compuestos}
\label{181224}
\entrada{Divisores_compuestos}

\chapter{Número de divisores compuestos}
\label{181225}
\entrada{Numero_de_divisores_compuestos}

\chapter{Tablas de operaciones binarias}
\label{181226}
\entrada{Tablas_de_operaciones_binarias}

\chapter{Reconocimiento de conmutatividad}
\label{181227}
\entrada{Reconocimiento_de_conmutatividad}

\chapter{Árbol de subconjuntos}
\label{181228}
\entrada{Arbol_de_subconjuntos}

% ---------------------------------------------------------------------

\chapter{El teorema de Navidad de Fermat}
\label{181231}
\entrada{El_teorema_de_Navidad_de_Fermat}

\chapter{El 2019 es apocalíptico}
\label{190101}
\entrada{El_2019_es_apocaliptico}

\chapter{El 2019 es malvado}
\label{190102}
\entrada{El_2019_es_malvado}

\chapter{El 2019 es semiprimo}
\label{190103}
\entrada{El_2019_es_semiprimo}

\chapter{El 2019 es un número de la suerte}
\label{190104}
\entrada{El_2019_es_un_numero_de_la_suerte}

% ---------------------------------------------------------------------

\chapter{Cadena descendiente de subnúmeros}
\label{190107}
\entrada{Cadena_descendiente_de_subnumeros}

\chapter{Mínimo producto escalar}
\label{190108}
\entrada{Minimo_producto_escalar}

\chapter{Numeración de ternas de naturales}
\label{190109}
\entrada{Numeracion_de_ternas}

\chapter{Subárboles monovalorados}
\label{190110}
\entrada{Subarboles_monovalorados}

\chapter{Mayor prefijo con suma acotada}
\label{190111}
\entrada{Mayor_prefijo_con_suma_acotada}

% ---------------------------------------------------------------------

\chapter{Ofertas 3 por 2}
\label{190114}
\entrada{Ofertas_3_por_2}

\chapter{Representación de conjuntos mediante intervalos}
\label{190115}
\entrada{Representacion_de_conjuntos_mediante_intervalos}

\chapter{Números altamente compuestos}
\label{190116}
\entrada{Numeros_altamente_compuestos}

\chapter{Posiciones del 2019 en el número pi}
\label{190117}
\entrada{Posiciones_del_2019_en_el_numero_pi}

\chapter{Mínimo número de operaciones para transformar
  un número en otro}
\label{190118}
\entrada{Minimo_numero_de_operaciones_para_transformar_un_numero_en_otro}

%-------------------------------------------------------------------------------

\chapter{Intersección de listas infinitas crecientes}
\label{190121}
\entrada{Interseccion_de_listas_infinitas_crecientes}

\chapter{Soluciones de \(x^2 = y^3 = k\)}
\label{190122}
\entrada{Soluciones_de_x2y3k}

\chapter{Sucesión triangular}
\label{190123}
\entrada{Sucesion_triangular}

\chapter{Números primos en pi}
\label{190124}
\entrada{Numeros_primos_en_pi}

\chapter{Recorrido de árboles en espiral}
\label{190125}
\entrada{Recorrido_de_arboles_en_espiral}

%-------------------------------------------------------------------------------

\chapter{Números con dígitos 1 y 2}
\label{190128}
\entrada{Numeros_con_digitos_1_y_2}

\chapter{Árboles con n elementos}
\label{190129}
\entrada{Arboles_con_n_elementos}



\end{document}

%%% Local Variables: 
%%% mode: latex
%%% TeX-master: t
%%% End: 

